
\section{EM Wave Fundamentals}
\label{sec:Em:em_wave_fundementals}

\subsection{Plane Electromagnetic Waves}

A simple solution to the electromagnetic wave equation
is a plane wave propagating in the $+x$ direction,

\begin{equation}
\boxed{
\mathbf{E}(x,t) =
\operatorname{Re}\left[
\mathbf{E}_0 e^{i(kx-\omega t)}
\right]
}
\label{eq:em_plane_wave}
\end{equation}

where $\mathbf{E}_0$ is the complex electric field
amplitude, $k$ is the angular wavenumber, and
$\omega$ is the angular frequency.

The angular wavenumber is defined as

\begin{equation}
\boxed{k = \frac{2\pi}{\lambda}}
\label{eq:angular_wavenumber}
\end{equation}

where $\lambda$ is the wavelength. Physically, $k$
describes how rapidly the phase changes with
position, measured in radians per metre.

Similarly, the angular frequency is

\begin{equation}
\boxed{\omega = 2\pi f = \frac{2\pi}{T}}
\label{eq:angular_frequency}
\end{equation}

where $f$ is the frequency and $T$ is the period.
Physically, $\omega$ describes how rapidly the
phase changes with time, measured in radians
per second.

In three dimensions, the plane wave becomes

\begin{equation}
\mathbf{E}(\mathbf{r},t) =
\operatorname{Re}\left[
\mathbf{E}_0
e^{i(\mathbf{k}\cdot\mathbf{r}-\omega t)}
\right],
\label{eq:em_plane_wave_3d}
\end{equation}

where $\mathbf{k}$ is the wavevector, pointing
in the direction of propagation, with
magnitude $|\mathbf{k}|=k$.


\subsection{Relationship Between Electric and Magnetic Fields}


\subsubsection{Transverse Nature of EM Waves}

We can show that electromagnetic waves are transverse
using Gauss's laws in free space,

\begin{equation}
\boxed{\nabla\cdot\mathbf{E} = 0,
\qquad
\nabla\cdot\mathbf{B} = 0.}
\end{equation}

Consider a plane wave of the form

\begin{equation}
\mathbf{E}(\mathbf{r},t)
= \operatorname{Re}\left[
\mathbf{E}_0 e^{i(\mathbf{k}\cdot\mathbf{r}-\omega t)}
\right],
\end{equation}

where $\mathbf{k}$ is the wavevector and
$\mathbf{E}_0$ is the constant electric field amplitude.

Taking the divergence of the complex plane-wave
representation gives

\begin{equation}
\nabla\cdot\left(
\mathbf{E}_0 e^{i(\mathbf{k}\cdot\mathbf{r}-\omega t)}
\right)
=
i(\mathbf{k}\cdot\mathbf{E}_0)
e^{i(\mathbf{k}\cdot\mathbf{r}-\omega t)}.
\end{equation}

From Gauss's law, this must equal zero in free space.
Since the exponential is never zero, we obtain

\begin{equation}
\boxed{\mathbf{k}\cdot\mathbf{E}_0 = 0}.
\end{equation}

Repeating the same calculation for the magnetic field,
using $\nabla\cdot\mathbf{B}=0$, gives

\begin{equation}
\boxed{\mathbf{k}\cdot\mathbf{B}_0 = 0}.
\end{equation}

Therefore,

\begin{equation}
\boxed{
\mathbf{k}\cdot\mathbf{E}_0 = 0,
\qquad
\mathbf{k}\cdot\mathbf{B}_0 = 0
}
\label{eq:em_transverse}
\end{equation}

Both the electric and magnetic fields are perpendicular
to the direction of wave propagation, $\mathbf{k}$.
This is why electromagnetic plane waves in free space
are called \textbf{transverse waves}.


and Faraday's law gives

\begin{equation}
\boxed{
\mathbf{B}_0 =
\frac{1}{\omega}
\mathbf{k}\times\mathbf{E}_0
}
\label{eq:em_e_b_relation}
\end{equation}

Thus, the electric field, magnetic field,
and propagation direction are mutually
perpendicular.

In free space, their amplitudes satisfy

\begin{equation}
\boxed{E_0 = c B_0}
\label{eq:em_field_amplitudes}
\end{equation}

where $c$ is the speed of light.


\subsection{Phase and Group Velocity}

The phase of a plane electromagnetic wave is the argument of the complex exponetial.

\begin{equation}
\boxed{\phi = kx-\omega t}.
\end{equation}

To follow a point of constant phase, we set
$d\phi=0$, giving

\begin{equation}
k\,dx-\omega\,dt=0.
\end{equation}

Therefore, the phase velocity is

\begin{equation}
\boxed{
v_{\mathrm{ph}} =
\frac{\omega}{k}
}
\label{eq:phase_velocity}
\end{equation}

\textbf{The phase velocity describes the speed at
which a point of constant phase, such as a
wave crest, propagates.
}
A wave packet can be constructed by
superposing waves with slightly different
frequencies and wavenumbers. The speed of
the resulting envelope is called the
\textbf{group velocity},

\begin{equation}
\boxed{
v_{\mathrm{g}} =
\frac{d\omega}{dk}
}
\label{eq:group_velocity}
\end{equation}

The group velocity describes the propagation
of the wave packet envelope. In a transparent,
weakly absorbing medium it is often associated
with the transport of energy.

In free space, the dispersion relation is

\begin{equation}
\boxed{\omega = ck}
\label{eq:vacuum_dispersion_relation}
\end{equation}

and consequently,

\begin{equation}
\boxed{
v_{\mathrm{ph}} = v_{\mathrm{g}} = c
}
\label{eq:vacuum_velocities}
\end{equation}


\subsection{Dispersion}

Dispersion occurs when the relationship
between angular frequency $\omega$ and
wavenumber $k$ is not linear.

In a dispersive medium, waves of different
frequencies can propagate with different
phase and group velocities.

\subsubsection{Dispersion in Free Space}

In classical vacuum electrodynamics,
the dispersion relation is

\begin{equation}
\omega^2 = c^2 k^2.
\end{equation}

Because $\omega=ck$ for positive-frequency
waves, both phase and group velocities are
independent of frequency.

Therefore, free space is a non-dispersive
medium.


\subsubsection{Electromagnetic Waves in Matter}

Maxwell's equations remain valid in matter,
but the electromagnetic response of the
material must be included.

For macroscopic electromagnetic fields,
we introduce the electric displacement
field $\mathbf{D}$ and magnetic field
intensity $\mathbf{H}$,

\begin{equation}
\mathbf{D} = \epsilon\mathbf{E},
\qquad
\mathbf{B} = \mu\mathbf{H},
\label{eq:material_constitutive}
\end{equation}

where $\epsilon$ and $\mu$ are the permittivity
and permeability of a linear, isotropic medium.

In a homogeneous, source-free, lossless,
non-dispersive medium, Maxwell's equations
give the wave equation

\begin{equation}
\boxed{
\nabla^2\mathbf{E} =
\mu\epsilon
\pdv[2]{\mathbf{E}}{t}
}
\label{eq:em_wave_equation_medium}
\end{equation}

The wave propagation speed is therefore

\begin{equation}
\boxed{
v = \frac{1}{\sqrt{\mu\epsilon}}
}
\label{eq:em_speed_medium}
\end{equation}

The refractive index is defined as

\begin{equation}
\boxed{
n = \frac{c}{v}
}
\label{eq:refractive_index}
\end{equation}

For a frequency-dependent material response,
the refractive index can depend on frequency,

\begin{equation}
n=n(\omega).
\end{equation}

For a transparent, isotropic medium, the
dispersion relation is then

\begin{equation}
\boxed{
k(\omega) = \frac{n(\omega)\omega}{c}
}
\label{eq:material_dispersion_relation}
\end{equation}

The phase velocity becomes

\begin{equation}
\boxed{
v_{\mathrm{ph}} = \frac{c}{n(\omega)}
}
\label{eq:material_phase_velocity}
\end{equation}

while the group velocity is

\begin{equation}
\boxed{
v_{\mathrm{g}} =
\frac{c}{
n(\omega)+\omega\dv{n}{\omega}
}
}
\label{eq:material_group_velocity}
\end{equation}

Thus, when the refractive index depends on
frequency, the phase and group velocities
are generally different.


% \subsubsection{Dispersion in a Cold Plasma}

% An important example in radio astronomy is
% electromagnetic wave propagation through
% an unmagnetized, cold electron plasma.

% The plasma frequency is

% \begin{equation}
% \boxed{
% \omega_{\mathrm{p}} =
% \sqrt{\frac{n_e e^2}{\epsilon_0 m_e}}
% }
% \label{eq:plasma_frequency}
% \end{equation}

% where $n_e$ is the electron number density,
% $e$ is the elementary charge, and $m_e$
% is the electron mass.

% The plasma dispersion relation is

% \begin{equation}
% \boxed{
% \omega^2 = \omega_{\mathrm{p}}^2+c^2k^2
% }
% \label{eq:plasma_dispersion_relation}
% \end{equation}

% The corresponding refractive index is

% \begin{equation}
% n(\omega) =
% \sqrt{1-\frac{\omega_{\mathrm{p}}^2}{\omega^2}}.
% \label{eq:plasma_refractive_index}
% \end{equation}

% For propagating waves with
% $\omega>\omega_{\mathrm{p}}$, the phase
% and group velocities are

% \begin{equation}
% v_{\mathrm{ph}} =
% \frac{c}{
% \sqrt{1-\omega_{\mathrm{p}}^2/\omega^2}
% },
% \end{equation}

% \begin{equation}
% \boxed{
% v_{\mathrm{g}} =
% c\sqrt{1-\frac{\omega_{\mathrm{p}}^2}{\omega^2}}
% }
% \label{eq:plasma_group_velocity}
% \end{equation}

% Notice that $v_{\mathrm{ph}}>c$, while
% $v_{\mathrm{g}}<c$. A superluminal phase
% velocity does not imply superluminal
% information transfer.

% Lower-frequency radio waves have smaller
% group velocities and therefore arrive
% later than higher-frequency waves.

% This frequency-dependent propagation delay
% is the origin of the dispersion observed
% in pulsars and fast radio bursts.
